\mode*
\section{Introduction}

Debugging is a substantial part of programming: things rarely work on the
first attempt, and the programmer must find out why.
Yet debugging is seldom taught explicitly in introductory programming courses;
students are largely left to pick it up on their own
\autocite[For a review of debugging in computing education, see][]{DebuggingReview}.
If we want to teach debugging, we first need a theory of what debugging
\emph{is}---what the skilled debugger does that the novice does not.

\mode<presentation>{%
\begin{frame}
  \begin{block}{Debugging}
    \begin{itemize}
      \item A substantial part of programming.
      \item Seldom taught explicitly.
      \item What \emph{is} the skill of debugging?
    \end{itemize}
  \end{block}
\end{frame}%
}

In this paper we propose such a theory, grounded in variation theory
\autocite{NCOL}.
The argument runs as follows.
A bug means that the program behaves differently from what the programmer
expects: the programmer has not yet discerned some critical aspect of the
program.
Debugging is therefore an act of \emph{learning}.
According to variation theory there is no learning without discernment, and no
discernment without variation; but during debugging there is no teacher to
provide that variation---the debugger must introduce it themselves.
\Textcite[Ch.~5]{NCOL} accounts for exactly this kind of teacher-less learning
in two places: in scientific discovery, where the researcher opens up new dimensions
of variation, and in the deep approach to learning
\autocite{MartonSaljo1976I}, where
\textcquote[p.~145]{NCOL}{the learner has to create the necessary patterns of
variation and invariance}.
This suggests a hypothesis: \emph{debugging skill and the deep approach to
learning are the same underlying ability}---the ability to generate the
appropriate patterns of variation for oneself.

\mode<presentation>{%
\begin{frame}
  \begin{block}{The argument}
    \begin{itemize}
      \item A bug: the program behaves differently than the programmer
        expects.
      \item The programmer hasn't discerned some critical aspect.
      \item Debugging is learning---without a teacher.
      \item The debugger must generate the variation themselves.
    \end{itemize}
  \end{block}

  \pause

  \begin{idea}
    \begin{itemize}
      \item Hypothesis: debugging skill \(=\) deep approach to learning
        \(=\) generating the right patterns of variation for oneself.
    \end{itemize}
  \end{idea}
\end{frame}%
}

This paper is the full version of the debugging chapter in our companion paper
on misconceptions in introductory programming
\autocite{VTProgMisconceptions}.
We have observed, informally, that experienced programmers who debug an
unfamiliar phenomenon generate exactly the patterns of variation that
variation theory prescribes for teaching it; the companion paper only
sketches this observation, here we report it in detail (\cref{pilot}).
We develop the observation into a testable hypothesis and outline a
study to test it: in the 2026 instance of an introductory programming course
(prgi26), we record the students' every edit--run cycle with the
\texttt{learnlog} tool \autocite{learnlog}, measure their approaches to
learning with the R-SPQ-2F questionnaire \autocite{RSPQ2F}, and analyse the
recorded debugging episodes for patterns of variation.
We also want an expert baseline to compare the students' patterns against;
as an approximation of expert debuggers, we record and analyse the debugging
of the course's teachers and teaching assistants in the same way.

\subsection{Research questions}

The purpose of this study is to answer the following questions:
\begin{frame}[allowframebreaks]
\begin{restatable}{question}{rqpatterns}\label{rq:patterns}
  Which patterns of variation (contrast, generalisation, fusion), if any,
  can be observed in students' recorded debugging episodes?
\end{restatable}
\begin{restatable}{question}{rqdeep}\label{rq:deep}
  Is a deep approach to learning associated with generating more, and more
  complete, patterns of variation while debugging?
\end{restatable}
\begin{restatable}{question}{rqsuccess}\label{rq:success}
  Is debugging success associated with the deep approach and with the
  completeness of the generated patterns?
\end{restatable}
\begin{restatable}{question}{rqexperts}\label{rq:experts}
  Which patterns of variation do expert debuggers exhibit while debugging,
  and how do the students' patterns compare to theirs?
\end{restatable}
\begin{restatable}{question}{rqgenerative}\label{rq:generative}
  Do students who have experienced the course material's patterns of
  variation generate more, and more complete, patterns of variation in
  their debugging than students who study on their own?
\end{restatable}
\end{frame}

These questions call for different types of methods.
\Cref{rq:patterns,rq:experts} are qualitative: they require recording actual
debugging episodes---of students and of expert debuggers, the latter
approximated by the course's teachers---and coding the recorded acts, what
was varied and what was kept invariant, in variation-theoretic terms.
In phenomenographic terms, we study the acts in which the debugger's
understanding is expressed.
\Cref{rq:deep,rq:success} are correlational: they require independent,
quantitative measures---of the approach to learning (the questionnaire), of
the generated patterns (the coded episodes) and of debugging success (the
outcome measures)---collected from the same students.
\Cref{rq:generative} anticipates the generative-learning corollary we
derive in \cref{debugging-as-learning}, where we return to it; it is
correlational as well, comparing groups within the cohort.
Together this makes for a mixed-methods design, which we detail in
\cref{method}.

The rest of the paper is structured as follows.
In \cref{background} we recap phenomenography, variation theory and the
approaches to learning.
In \cref{debugging-as-learning} we develop the theoretical argument that
debugging is self-directed learning and state our hypotheses.
In \cref{debugging} we give a worked example of the patterns of variation in
an actual debugging episode.
In \cref{method} we describe the planned study and in \cref{results} the
expected results.
In \cref{related-work} we compare with related work---placed late so that
the comparison with our design is direct---and in \cref{discussion} we
conclude.
